\subsubsection{Arithmetization over \texorpdfstring{$\FF_2$}{F2}}
\label{s:sha256-f2-arithmetization}

We now give a self-contained arithmetization of the complete SHA-256
compression function whose witness symbols, coefficients, and constraints all
belong to $\FF_2$.  The construction follows the word-level decomposition of
Zinc+~\cite[Section~8.2]{zincplus}, but it replaces the integer modular-sum
identity used there by a Boolean carry transition.  This replacement is
necessary: reducing an identity of the form
\[
  \sum_j x_j+c_{\mathrm{in}}
  =y+2c_{\mathrm{out}}
\]
modulo $2$ deletes the term containing $c_{\mathrm{out}}$ and therefore does
not constrain the carry.

\paragraph{Words and the linear layer.}
We identify a word $x$ with its little-endian bit vector
$x=(x_0,\ldots,x_{31})\in\FF_2^{32}$ and write
\[
  \langle x\rangle_2\coloneqq\sum_{i=0}^{31}2^i x_i\in
  \{0,\ldots,2^{32}-1\}.
\]
For $0\leq r<32$, define the public $\FF_2$-linear maps
\[
  (R_r x)_i=x_{(i+r)\bmod 32},
  \qquad
  (S_r x)_i=
  \begin{cases}
    x_{i+r},&i+r<32,\\
    0,&i+r\geq32.
  \end{cases}
\]
Thus $R_r$ and $S_r$ are respectively $\ROTR^r$ and $\SHR^r$ under the
little-endian convention.  Consequently,
\begin{equation}\label{eq:sha-f2-linear-layer}
\begin{aligned}
\Sigma_0(x)&=R_2x+R_{13}x+R_{22}x,&
\Sigma_1(x)&=R_6x+R_{11}x+R_{25}x,\\
\sigma_0(x)&=R_7x+R_{18}x+S_3x,&
\sigma_1(x)&=R_{17}x+R_{19}x+S_{10}x
\end{aligned}
\end{equation}
are linear maps over $\FF_2$.  Every addition in
\eqref{eq:sha-f2-linear-layer} is XOR.  These values may be materialized as
linear witness columns, or exposed virtually from the committed columns by the
$\FF_2$-linear virtualization mechanism of \cref{s:virtualization}.

\paragraph{Choice and majority.}
For bits $x,y,z\in\FF_2$, Boolean AND is multiplication, NOT is addition by
$1$, and XOR is addition.  The two nonlinear SHA-256 gates therefore have the
following rank-one forms
\begin{equation}\label{eq:sha-f2-ch-maj-anf}
  \Ch(x,y,z)=z+x(y+z),
  \qquad
  \Maj(x,y,z)=x+(x+y)(x+z).
\end{equation}
Each expression uses one multiplication.  The majority identity follows by
expanding and using $x^2=x$ to obtain $xy+xz+yz$; this Boolean identity is
automatic for an $\FF_2$-typed wire.  For word inputs, the
equations are imposed coordinate-wise.  Concretely, in round $t$ we commit
two product words $u^{\Ch}_t,u^{\Maj}_t$ and impose, for every
$i\in\{0,\ldots,31\}$,
\begin{equation}\label{eq:sha-f2-gate-constraints}
\begin{aligned}
 u^{\Ch}_{t,i}      &=e_{t,i}(f_{t,i}+g_{t,i}),
 &\Ch_{t,i}         &=g_{t,i}+u^{\Ch}_{t,i},\\
 u^{\Maj}_{t,i}     &=(a_{t,i}+b_{t,i})(a_{t,i}+c_{t,i}),
 &\Maj_{t,i}         &=a_{t,i}+u^{\Maj}_{t,i}.
\end{aligned}
\end{equation}
Only the two product equations are quadratic; the displayed outputs are
linear combinations and can again be virtual.  These are coordinate-wise
(Hadamard) products of bit columns, not convolution products in
$\FF_2[X]$.

\paragraph{A native characteristic-two adder.}
We next give an exact $k$-operand adder over $\FF_2$.  This is the only part of
the construction for which merely rewriting the usual integer identity in
characteristic two is unsound.

For a finite list of bits $U=(u_1,\ldots,u_L)$ and $d\geq0$, let
\[
  E_d(U)\coloneqq
  \sum_{\substack{A\subseteq[L]\\|A|=d}}\ \prod_{j\in A}u_j
  \quad\in\FF_2
\]
be its degree-$d$ elementary symmetric polynomial, with $E_0(U)=1$.
Repeated entries of $U$ are treated as distinct positions.  Since the
variables are bits, we always use the multilinear reduction of this
polynomial.

\begin{lemma}[Binary digit extraction]\label{lem:sha-f2-digit-extraction}
Let $U\in\FF_2^L$, let $H=\sum_{j=1}^L U_j$ be its Hamming weight as an
integer, and write $H=\sum_{q\geq0}2^qH_q$.  Then
\[
  E_{2^q}(U)=H_q\qquad\text{in }\FF_2.
\]
\end{lemma}
\begin{proof}
If $H$ entries of $U$ are one, then
$E_d(U)=\binom{H}{d}\bmod2$.  Lucas' theorem applied to
$d=2^q$ gives $\binom{H}{2^q}\bmod2=H_q$.
\end{proof}

Fix $k\geq2$ and put $\ell_k=\lceil\log_2 k\rceil$.  To enforce
\[
  y=\sum_{j=1}^k x^{(j)}\pmod{2^{32}},
\]
introduce carry bits
$\gamma_{i,p}\in\FF_2$ for
$i\in\{0,\ldots,32\}$ and $p\in\{0,\ldots,\ell_k-1\}$.
They encode the integer
$c_i=\sum_p2^p\gamma_{i,p}$ entering bit position $i$.  Set
$\gamma_{0,p}=0$.  At bit $i$, form the list
\begin{equation}\label{eq:sha-f2-carry-multiset}
  U_i^{(k)}
  \coloneqq
  \bigl(
    x^{(1)}_i,\ldots,x^{(k)}_i,\,
    \gamma_{i,0},\,
    \underbrace{\gamma_{i,1},\gamma_{i,1}}_{2\text{ copies}},\,
    \ldots,\,
    \underbrace{\gamma_{i,\ell_k-1},\ldots,
                \gamma_{i,\ell_k-1}}_{2^{\ell_k-1}\text{ copies}}
  \bigr).
\end{equation}
Its integer Hamming weight is
$\sum_jx_i^{(j)}+c_i$.  We impose
\begin{equation}\label{eq:sha-f2-add-k}
\boxed{
\begin{aligned}
 y_i+E_1\bigl(U_i^{(k)}\bigr)&=0,\\
 \gamma_{i+1,p}
   +E_{2^{p+1}}\bigl(U_i^{(k)}\bigr)&=0
       &&(0\leq p<\ell_k)
\end{aligned}}
\qquad(0\leq i<32).
\end{equation}
The first equation is linear:
$E_1(U_i^{(k)})=\sum_jx_i^{(j)}+\gamma_{i,0}$, because every
higher carry digit occurs an even number of times.

\begin{lemma}[Exactness of the $\FF_2$ adder]\label{lem:sha-f2-adder}
With $\gamma_0=0$, \eqref{eq:sha-f2-add-k} holds if and only if
\[
  \langle y\rangle_2
  =\sum_{j=1}^k\langle x^{(j)}\rangle_2\pmod{2^{32}},
\]
and every $\gamma_i$ is the binary encoding of the true carry entering
position $i$.  In particular, the carry is uniquely determined.
\end{lemma}
\begin{proof}
Assume the carry entering position $i$ is correct.  By
\cref{lem:sha-f2-digit-extraction}, the first equation of
\eqref{eq:sha-f2-add-k} selects bit zero of
$c_i+\sum_jx_i^{(j)}$, while its remaining equations select bits
$1,\ldots,\ell_k$ of the same integer.  These are respectively the output
bit $y_i$ and the binary digits of
$c_{i+1}=\lfloor(c_i+\sum_jx_i^{(j)})/2\rfloor$.
The claim follows by induction from $c_0=0$.  The last carry $c_{32}$ is the
discarded overflow, which is exactly reduction modulo $2^{32}$.
\end{proof}

The $\FF_2$ typing is load-bearing.  If these equations are evaluated inside
an extension field, every operand, output, product, and carry wire must still
be constrained to the prime subfield (equivalently, satisfy $v^2=v$).
Without that condition, the Hamming-weight argument in
\cref{lem:sha-f2-digit-extraction} does not apply.

For $k=2$, \eqref{eq:sha-f2-add-k} is the familiar full adder
\begin{equation}\label{eq:sha-f2-full-adder}
  y_i=x_i^{(1)}+x_i^{(2)}+\gamma_i,
  \qquad
  \gamma_{i+1}
   =x_i^{(1)}x_i^{(2)}
    +x_i^{(1)}\gamma_i+x_i^{(2)}\gamma_i.
\end{equation}
For the arities needed by SHA-256, the compact encoding has the following
cost.  A ``carry plane'' consists of the $32$ carries produced by one word
addition; the final, discarded carry may instead be omitted if a packed
layout does not require full words.
\begin{table}[H]
\centering
\small
\caption{Native $\FF_2$ adders used by the SHA-256 arithmetization.}
\label{tab:sha-f2-adders}
\begin{tabular}{@{}cccccc@{}}
\toprule
Operands $k$ & Carry planes $\ell_k$ & Max.\ degree
  & Linear eqs. & Nonlinear eqs. & Transition-table rows\\
\midrule
$2$ & $1$ & $2$ & $32$ & $32$ & $8$\\
$3$ & $2$ & $4$ & $32$ & $64$ & $24$\\
$4$ & $2$ & $4$ & $32$ & $64$ & $64$\\
$5$ & $3$ & $7$ & $32$ & $96$ & $160$\\
\bottomrule
\end{tabular}
\end{table}
For $k=5$, the highest raw digit extractor is $E_8$; its multilinear
reduction on the eight distinct input and carry variables has degree seven,
which is the degree reported in \cref{tab:sha-f2-adders}.
For implementation, let $e_r$ be the degree-$r$ elementary symmetric
polynomial in the $k$ operand bits at one position.  In the Boolean quotient,
the required digit polynomials reduce to
\begin{equation}\label{eq:sha-f2-reduced-digits}
\begin{aligned}
E_1&=e_1+\gamma_0,\\
E_2&=e_2+\gamma_0e_1+\gamma_1,\\
E_4&=e_4+\gamma_0e_3+\gamma_1e_2
       +\gamma_0\gamma_1e_1+\gamma_2,\\
E_8&=e_8+\gamma_0e_7+\gamma_1e_6+\gamma_0\gamma_1e_5
       +\gamma_2e_4+\gamma_0\gamma_2e_3\\
   &\hspace{2.8em}+\gamma_1\gamma_2e_2
       +\gamma_0\gamma_1\gamma_2e_1.
\end{aligned}
\end{equation}
Terms containing an unavailable carry digit or $e_r$ with $r>k$ are omitted.
These formulas follow by extracting coefficients from
\[
 \prod_{j=1}^k(1+x_i^{(j)}T)
 \prod_{p=0}^{\ell_k-1}(1+\gamma_{i,p}T^{2^p}).
\]

There are two useful implementation choices.  A backend accepting bounded
degree constraints can use \eqref{eq:sha-f2-add-k} directly.  A backend with
binary lookups can instead require each local tuple
\[
 \bigl(x_i^{(1)},\ldots,x_i^{(k)},\gamma_i,y_i,\gamma_{i+1}\bigr)
\]
to belong to the fixed transition table defined by ordinary binary addition.
The table has $k2^k$ rows, as shown in \cref{tab:sha-f2-adders}, and all of
its entries are in $\FF_2$.  Finally, a quadratic-only backend can replace a
$k$-input adder by $k-1$ chained copies of
\eqref{eq:sha-f2-full-adder}.  This last option lowers the degree at the cost
of additional carry planes and greater multiplicative depth.

\paragraph{Composing the compression function.}
Write
\[
 \mathsf{Add}_k(x^{(1)},\ldots,x^{(k)};y,\gamma)
\]
for the constraints in \eqref{eq:sha-f2-add-k}.  All arguments are
$32$-bit vectors over $\FF_2$.  The complete SHA-256 compression
arithmetization is as follows.

\begin{enumerate}[label=\textbf{F\arabic*.},leftmargin=2.8em]
\item\label{item:sha-f2-input}
  The witness contains the input chaining state
  $H^{\mathrm{in}}\in(\FF_2^{32})^8$, the sixty-four schedule words
  $W_1,\ldots,W_{64}$, and the round values specified below.  Bind the first
  sixteen schedule words to the provided message words and the initial
  registers to the input state:
  \begin{equation}\label{eq:sha-f2-boundary}
  W_t=M_t\quad(t\in[16]),
  \qquad
  (a_1,b_1,c_1,d_1,e_1,f_1,g_1,h_1)=H^{\mathrm{in}}.
  \end{equation}

\item\label{item:sha-f2-schedule}
  For each $t\in\{17,\ldots,64\}$, expose
  $\sigma_0(W_{t-15})$ and $\sigma_1(W_{t-2})$ by
  \eqref{eq:sha-f2-linear-layer}, and impose
  \begin{equation}\label{eq:sha-f2-schedule}
  \mathsf{Add}_4\bigl(
       W_{t-16},\sigma_0(W_{t-15}),W_{t-7},\sigma_1(W_{t-2});
       W_t,\gamma^W_t
  \bigr).
  \end{equation}

\item\label{item:sha-f2-rounds}
  For every $t\in[64]$, enforce the two product words in
  \eqref{eq:sha-f2-gate-constraints}, expose the two big-$\Sigma$ values and
  the required XORs linearly, and impose
  \begin{equation}\label{eq:sha-f2-round-adders}
  \begin{aligned}
  \mathsf{Add}_5\bigl(
    h_t,\Sigma_1(e_t),\Ch(e_t,f_t,g_t),K_t,W_t;
    T_{1,t},\gamma^{T_1}_t\bigr),\\
  \mathsf{Add}_3\bigl(
    T_{1,t},\Sigma_0(a_t),\Maj(a_t,b_t,c_t);
    a_{t+1},\gamma^a_t\bigr),\\
  \mathsf{Add}_2\bigl(
    d_t,T_{1,t};e_{t+1},\gamma^e_t\bigr).
  \end{aligned}
  \end{equation}
  The remaining registers are wiring:
  \begin{equation}\label{eq:sha-f2-register-shift}
  (b_{t+1},c_{t+1},d_{t+1},f_{t+1},g_{t+1},h_{t+1})
  =(a_t,b_t,c_t,e_t,f_t,g_t).
  \end{equation}
  Notice that \eqref{eq:sha-f2-round-adders} absorbs
  $T_{2,t}=\Sigma_0(a_t)+\Maj(a_t,b_t,c_t)$ directly into the
  three-input addition for $a_{t+1}$, so no $T_2$ word is committed.

\item\label{item:sha-f2-feed-forward}
  Let
  $S^{\mathrm{fin}}=(a_{65},b_{65},c_{65},d_{65},
  e_{65},f_{65},g_{65},h_{65})$.  For each $j\in\{0,\ldots,7\}$, impose
  \begin{equation}\label{eq:sha-f2-feed-forward}
    \mathsf{Add}_2\bigl(
      H^{\mathrm{in}}_j,S^{\mathrm{fin}}_j;
      H^{\mathrm{out}}_j,\gamma^{\mathrm{ff}}_j
    \bigr).
  \end{equation}
  The eight words $H^{\mathrm{out}}$ are the output chaining state.
\end{enumerate}

\begin{theorem}[Completeness and exactness]\label{thm:sha-f2-correctness}
The constraints
\eqref{eq:sha-f2-linear-layer},
\eqref{eq:sha-f2-gate-constraints},
\eqref{eq:sha-f2-boundary},
\eqref{eq:sha-f2-schedule},
\eqref{eq:sha-f2-round-adders},
\eqref{eq:sha-f2-register-shift}, and
\eqref{eq:sha-f2-feed-forward}
are satisfiable if and only if $H^{\mathrm{out}}$ is the SHA-256 compression
of the message block from initial chaining state $H^{\mathrm{in}}$.
For fixed input words, every schedule, state, product word, and carry is
uniquely determined.
\end{theorem}
\begin{proof}
The maps in \eqref{eq:sha-f2-linear-layer} are exactly the FIPS rotations,
shifts, and XORs.  Equation \eqref{eq:sha-f2-ch-maj-anf} proves the
choice and majority constraints coordinate-wise.  Applying
\cref{lem:sha-f2-adder} to \eqref{eq:sha-f2-schedule} gives the standard
message recurrence modulo $2^{32}$.  The same lemma applied to
\eqref{eq:sha-f2-round-adders}, followed by the wiring in
\eqref{eq:sha-f2-register-shift}, gives each of the sixty-four standard
round transitions.  A final application to
\eqref{eq:sha-f2-feed-forward} gives the feed-forward state.  This proves
soundness by induction over schedule indices, then round indices.
Conversely, a genuine execution supplies the unique products and carries and
therefore satisfies every equation.
\end{proof}

\paragraph{Proof-obligation ledger.}
\Cref{tab:sha-f2-ledger} records the reduction underlying the preceding
proof.  Each row shows what remains to be justified after that stage; the
last row is empty.
\begin{table}[H]
\centering
\small
\caption{Claim ledger for the all-$\FF_2$ construction.}
\label{tab:sha-f2-ledger}
\begin{tabular}{@{}p{0.22\textwidth}p{0.29\textwidth}p{0.39\textwidth}@{}}
\toprule
Stage & Discharged here & Obligations still open\\
\midrule
Relation statement
  & \textcolor{ForestGreen}{Created:} correct compression
  & linear layer; Boolean gates; modular addition; composition\\
Linear layer
  & rotations, shifts, XOR
  & Boolean gates; modular addition; composition\\
\Cref{eq:sha-f2-gate-constraints}
  & choice and majority
  & modular addition; composition\\
\Cref{lem:sha-f2-adder}
  & every modular word sum
  & schedule, rounds, and feed-forward composition\\
\Cref{thm:sha-f2-correctness}
  & complete SHA-256 compression
  & \textcolor{BrickRed}{none}\\
\bottomrule
\end{tabular}
\end{table}

\paragraph{Concrete cost.}
One implementation-oriented layout uses exactly the values that are nonlinear
inputs or outputs as committed words and exposes every $\FF_2$-linear value
virtually.  Its committed witness per compression is
\begin{equation}\label{eq:sha-f2-witness-cost}
\underbrace{8}_{H^{\mathrm{in}}}
+\underbrace{64}_{W}
+\underbrace{48\cdot2}_{\text{schedule carry planes}}
+\underbrace{64\cdot11}_{\text{round words}}
+\underbrace{8}_{H^{\mathrm{out}}}
+\underbrace{8}_{\text{feed-forward carries}}
=888
\end{equation}
$32$-bit words, or $28{,}416$ committed bits.  The eleven words per round
are $T_1$, its three carry planes, $a_{t+1}$, its two carry planes,
$e_{t+1}$, its carry plane, and the two product words from
\eqref{eq:sha-f2-gate-constraints}.  Shifted registers reuse earlier $a$ and
$e$ words.  The virtual linear layer has
\[
  48\cdot2+64\cdot7=544
\]
words: two small-sigma words per derived schedule entry and, per round,
$\Sigma_0$, $\Sigma_1$, $f+g$, $\Ch$, $a+b$, $a+c$, and $\Maj$.

Using the symmetric-polynomial adders directly gives $7{,}936$ linear
sum-bit equations and the following nonlinear equations:
\begin{center}
\begin{tabular}{@{}l>{\raggedright\arraybackslash}p{0.62\textwidth}r@{}}
\toprule
Maximum degree & Source & Number\\
\midrule
$2$ & all $E_2$ carry digits ($7{,}936$) and the two product words per
      round ($4{,}096$)
    & $12{,}032$\\
$4$ & the $E_4$ digit of the $3$-, $4$-, and $5$-input adders
    & $5{,}632$\\
$7$ & the multilinear reduction of the $E_8$ digit of the $5$-input round adder
    & $2{,}048$\\
\midrule
Total & & $19{,}712$\\
\bottomrule
\end{tabular}
\end{center}
The $544$ virtual words would add $17{,}408$ linear equations if
materialized instead.  A quadratic-only chain uses $600$ two-input word
adders, hence $19{,}200$ quadratic carry equations in addition to the
$4{,}096$ gate equations, and increases the compact committed layout from
$888$ to $1{,}000$ words (assuming intermediate sum words remain virtual).
Thus the direct construction trades degree at most seven for fewer witness
bits and fewer local constraints; the transition-table variant retains the
same compact witness without using high-degree carry polynomials.

For compatibility with a layout that already uses the two-product identity
$\Maj(a,b,c)=ab+c(a+b)$, replace $u^{\Maj}$ by two product words.  This removes
the virtual word $a+c$ but adds one committed word and $32$ quadratic
constraints per round, giving $952$ committed words, $480$ virtual words, and
$21{,}760$ nonlinear equations.  The schedule, carries, and soundness proof
are unchanged.
